% !PRINTED 514
\weberpage{514}
%<<CONT
worin die Quadratwurzeln positiv zu nehmen sind. Sodann ist, weil $B$ gerade ist
\[
\alpha m'' \equiv -1 \ (\mathrm{mod}\ 4), \quad \left(\frac{\beta}{\alpha}\right)=\left(\frac{A}{\alpha}\right),
\]
und folglich nach dem Reziprozitätsgesetz und nach (4):
\[
\left(\frac{\alpha m''}{A}\right) = \left(\frac{-1}{A}\right) = \left(\frac{-1}{A}\right)\left(\frac{A}{\alpha m''}\right),
\]
also:
\[
\left(\frac{A}{\alpha m''}\right) = 1.
\]
\weq{6}{\left(\frac{\beta}{\alpha}\right) = \left(\frac{A}{m''}\right),}
\weq{7}{\beta \equiv 4, \quad \gamma \equiv 4\,C \ (\mathrm{mod}\ 8), \quad \alpha \equiv 1 \ (\mathrm{mod}\ 2) \quad [\S\ 136, (7)],}
also:
\[
e^{\frac{\pi i \alpha}{4}(\gamma-\beta)} = (-1)^{C-1},
\]
ferner:
\[
B + 2m' \equiv 2C \ (\mathrm{mod}\ 4),
\]
also:
\[
\frac{\alpha-1}{2} \equiv -\frac{m''+1}{2} + 2C \ (\mathrm{mod}\ 4),
\]
\[
i^{3\frac{\alpha-1}{2}} = (-1)^C i^{-3\frac{m''+1}{2}},
\]
und demnach endlich nach § 136, (12), (13):
\[
M = -\left(\frac{A}{m''}\right) i^{-3\frac{m''+1}{2}} \left(2\sqrt{m'} + i\sqrt{m''}\right)^3, % faint pre-minus speck is scan/print damage, not a character
\]
wofür man auch setzen kann:
\weq{8}{M = -\left(\frac{A}{m''}\right)\left(2\, i^{-\frac{m''+1}{2}}\sqrt{m'} + i^{-\frac{m''-1}{2}}\sqrt{m''}\right)^3 \, \cdot}
Nach dem Satz 1. (§ 136) ist dann $H_{-D}(u)$ teilbar durch eine Funktion
\[
\Phi(M,u),
\]
die zu $\Phi(-M,u)$ relativ prim ist. Setzen wir
\weq{9}{\mu = -\left(2\, i^{\frac{m''+1}{2}}\sqrt{m'} + i^{\frac{m''-1}{2}}\sqrt{m''}\right)^3,
}
so ist für jede Klasseninvariante von $D$ eine der beiden Gleichungen
\weq{10}{\Phi(\mu,u)=0, \quad \Phi(-\mu,u)=0}
%CONT>>
